\emph{(1).} Take generators $\gamma_{0}, \ldots, \gamma_{3}$ of spacetime and, for each rung plane $a = 4, 5$, $e_{a}$, $f_{a}$ with $e_{a}^{2} = 1$, $f_{a}^{2} = -1$, all mutually orthogonal, and $v_{a} = \tfrac12(e_{a} + f_{a})$, $w_{a} = \tfrac12(e_{a} - f_{a})$; then $v_{a}^{2} = w_{a}^{2} = 0$ and $\{v_{a}, w_{b}\} = \delta_{ab}$, the stored null pairs. A vector orthogonal to the four $\gamma_{\mu}$ anticommutes with each of them, so it commutes with their product $I_{4}$ (four sign changes); $I_{4}^{2} = -1$. This is the stored property that $I_{4}$ commutes with every mode (\claimref{obj.field-theory.single-algebra-field}), here for the rung modes. Hence $\mathrm{span}\{1, I_{4}\}$ acts on $H_{4} \oplus H_{5}$ as a complex structure inside the algebra, and for $c = \sum_{a}(\alpha_{a}v_{a} + \beta_{a}w_{a})$ with $\alpha_{a}, \beta_{a} \in \mathrm{span}\{1, I_{4}\}$, $c^{2} = \sum_{a,b}\alpha_{a}\beta_{b}\{v_{a}, w_{b}\} = \alpha^{\mathsf T}\beta = Q(c)$, bilinear over $\mathrm{span}\{1, I_{4}\}$: the stored complex-bilinear form, its unit read as $I_{4}$. (In the full algebra $I_{4}$ is not central: it anticommutes with the spacetime vectors, as the field theory uses; on the rung summand it commutes.) Four $\mathrm{span}\{1, I_{4}\}$ coefficients give four $I_{4}$-complex modes per momentum, each realised by the polynomial representation of \claimref{field-theory.fock-space-maxwell-modes}(1) with variables indexed by the rung field's modes.

Let $b_{a} = v_{a}w_{a} - \tfrac12 = v_{a}\wedge w_{a}$ and $b = b_{4} + b_{5}$. From $w_{a}v_{a} = 1 - v_{a}w_{a}$ and $v_{a}^{2} = 0$: $v_{a}w_{a}v_{a} = v_{a}$, so $[b, v_{a}] = v_{a}$, and likewise $[b, w_{a}] = -w_{a}$; $b_{4}$ and $b_{5}$ commute (disjoint planes), $b_{a}^{2} = v_{a}w_{a}v_{a}w_{a} - v_{a}w_{a} + \tfrac14 = \tfrac14$, and $(I_{4}b_{a})^{2} = -\tfrac14$. So $e^{tb_{a}} = \cosh\tfrac t2 + 2b_{a}\sinh\tfrac t2$, $e^{\phi I_{4}b_{a}} = \cos\tfrac\phi2 + 2I_{4}b_{a}\sin\tfrac\phi2$, and $\mathrm{Ad}(e^{zb}) = e^{z\,\mathrm{ad}_{b}}$ multiplies $v_{a}$ by $e^{z}$ and $w_{a}$ by $e^{-z}$, for $z = t$ real (the stored $e^{tK}$, $Kv_{a} = v_{a}$, $Kw_{a} = -w_{a}$, \claimref{synthesis.rung-field-on-the-lift-locus}(4)) and for $z = \phi I_{4}$. On coefficients the latter is $\alpha \mapsto e^{\phi I_{4}}\alpha$, $\beta \mapsto e^{-\phi I_{4}}\beta$, which is $e^{2\phi Y}$ for $Y = \tfrac{1}{2}I_{4}\mathbf{1}$ acting by $\alpha \mapsto Y\alpha$, $\beta \mapsto \bar{Y}\beta$. So $Y = \tfrac12 I_{4}K$ on the rung field: the unitary hypercharge flow is the rung boost at $I_{4}$-imaginary rapidity, which is the reading \claimref{field-theory.boost-modes} gives for any plane's modes (real rapidity scales, imaginary rapidity is the unitary phase). The identification is on rung values only: each $\gamma_{\mu}$ anticommutes with $I_{4}$ and commutes with $b$, so it anticommutes with $I_{4}b$, and $e^{\phi I_{4}b}\gamma_{\mu}e^{-\phi I_{4}b} = e^{2\phi I_{4}b}\gamma_{\mu} \ne \gamma_{\mu}$ for $\phi \notin \pi\mathbb{Z}$: the automorphism moves every spacetime vector and does not commute with the Lorentz rotors, so it is not a spacetime-point transformation, while on the rung summand it is the hypercharge phase. The check verifies all of this in $\Cl(3,5)$ at $e^{t} = 4$ and $(\cos\phi, \sin\phi) = (-\tfrac{7}{25}, \tfrac{24}{25})$, including that $\gamma_{0}$ anticommutes with $v_{a}$ and with $I_{4}b$ and that its image is not $\gamma_{0}$.

\emph{(2).} By \claimref{synthesis.rung-field-d-term-action-and-tree-masses-on-the-lift-locus}(1), on all of $\mathcal{R}_{\C}$, $\sum_{j}\mu_{j}^{2} = \tfrac14(N^{2} - |Q|^{2})$ and $\mu_{Y} = \tfrac12 m$, so $V = F(N, m) + (k_{10} - \tfrac18 g^{2})|Q|^{2} + k_{8}\operatorname{Re}Q^{2} + k_{9}\operatorname{Im}Q^{2}$ with $F = \tfrac18 g^{2}N^{2} + \tfrac12 g'^{2}(\tfrac12 m - \zeta)^{2} + W(N)$. Write $A = |\alpha|^{2}$, $B = |\beta|^{2}$, $N = \tfrac12(A + B)$, $m = \tfrac12(A - B)$. At $\alpha = 0$, $\beta = (-\lambda, 0)$: $\partial_{B}V = \tfrac12(F_{N} - F_{m})$, which vanishes by the stored vacuum condition (on the cone it reads $(g^{2} + g'^{2})\lambda^{2} + 4g'^{2}\zeta + 8W' = 0$), so $F_{m} = F_{N} = \tfrac14 g^{2}N^{*} + W'$. The $\alpha$ directions enter $V$ through $A$ (quadratically in $\alpha$) and through $Q = -\lambda\alpha_{1}$ (linearly in $\alpha_{1}$); hence the first derivatives in $\alpha$ vanish, the $\alpha$-$\beta$ second derivatives vanish, $\alpha_{2}$ does not mix with $\alpha_{1}$, and in real coordinates the Hessian on $\alpha_{2}$ is $2\partial_{A}V = F_{N} + F_{m} = \tfrac12 g^{2}N^{*} + 2W' = M_{W}^{2} + 2W'$, the stored charged mass, while on $\alpha_{1}$ it is $M_{W}^{2} + 2W'$ plus the Hessian of $\lambda^{2}\bigl((k_{10} - \tfrac18 g^{2})|\alpha_{1}|^{2} + k_{8}\operatorname{Re}\alpha_{1}^{2} + k_{9}\operatorname{Im}\alpha_{1}^{2}\bigr)$, the matrix $2\lambda^{2}\bigl((k_{10} - \tfrac18 g^{2})\mathbf{1} + \bigl(\begin{smallmatrix}k_{8} & k_{9}\\ k_{9} & -k_{8}\end{smallmatrix}\bigr)\bigr)$, with eigenvalues $4N^{*}(k_{10} \pm r) - M_{W}^{2}$. The sum is $2W' + 4N^{*}(k_{10} \pm r)$. The kinetic coefficient of each real coordinate is one ($2N(e) = 1$). The $\beta$ block is the stored one: three directions of the gauge orbit (rank three, the generator $Y - T_{3}$ fixing the point) are massless by invariance, and the radial direction has $m_{h}^{2} = M_{Z}^{2} + 2N^{*}W''$ (\claimref{synthesis.rung-field-d-term-action-and-tree-masses-on-the-lift-locus}(4) at $m^{*} = -N^{*}$, $H = 0$). Since $Q = \alpha^{\mathsf T}\beta$ has differential $\beta^{\mathsf T}d\alpha + \alpha^{\mathsf T}d\beta = -\lambda\,d\alpha_{1}$ at the point, the normal pair is exactly the pair of directions transverse to $L_{0}$. For the pure D-term ($W = 0$, $k = 0$) its mass vanishes: these are D-flat directions.

\emph{The level.} With $k_{8} = k_{9} = 0$ and the term $k_{10}|Q + H^{2}|^{2}$, restrict to $\alpha = (a, 0)$, $\beta = (-\lambda', 0)$ with $a, \lambda'$ real. Since $V$ depends on $a$ only through $a^{2}$ and through $Q = -a\lambda'$, the mixed second derivative in $(a, \lambda')$ vanishes at $a = 0$, so the shift of $\lambda'$ at first order in $H^{2}$ does not enter the $a$-equation at that order. Along real $\alpha_{1} = a$ the part of $V$ without the $k_{10}$ term is quadratic in $a$ at first order, with coefficient half the normal Hessian of (2) at $k_{10} = 0$, that is $W'$; so the $a$-derivative is $2W'a - 2\lambda k_{10}(H^{2} - a\lambda) + O(H^{4})$, which vanishes for $a(2W' + 4N^{*}k_{10}) = 2\lambda k_{10}H^{2}$, that is $a = \rho H^{2}/\lambda$ with $\rho = 1 - 2W'/m_{H}^{2}$, $m_{H}^{2} = 2W' + 4N^{*}k_{10}$, and $Q = -\rho H^{2}$. The check verifies this as the vanishing of the constant and linear coefficients in $H^{2}$ of the exact $\alpha_{1}$-gradient, and that $\rho = 1$ fails when $W' \ne 0$. As $k_{10} \to \infty$ with $m_{h}$, $m_{\pm}$ fixed, $\rho \to 1$ and the normal pair decouples: the constrained field is that limit. For $H \ne 0$ the off-diagonal Hessian entries between the normal and neutral directions are of relative order $H^{2}/N^{*}$ (computed symbolically), about $10^{-88}$ at the weak scale.

\emph{(3).} On the ray $\alpha = 0$, $V = k_{2}|\beta|^{4}$, so the ray is flat exactly when $k_{2} = 0$, which is $a + b - e = 2k_{2} = 0$ in the stored coordinates $a = k_{1} + k_{2} + k_{3}$, $b = k_{1} + k_{2} - k_{3}$, $e = 2(k_{1} - k_{2})$; then $e = 2k_{1}$ and $2a - e = 2k_{3}$. The only term linear in $\alpha$ is $|\beta|^{2}(k_{6}\operatorname{Re}Q + k_{7}\operatorname{Im}Q) = -\lambda^{3}(k_{6}\operatorname{Re}\alpha_{1} + k_{7}\operatorname{Im}\alpha_{1})$, a tadpole unless $k_{6} = k_{7} = 0$; the $k_{4}, k_{5}$ terms are cubic in $\alpha$ and do not enter the Hessian. As in (2), $\partial_{A}V = k_{3}\lambda^{2}$ gives $m_{\pm}^{2} = 2k_{3}\lambda^{2} = 2(2a - e)N$ on $\alpha_{2}$, and $\alpha_{1}$ adds $4N(k_{10} \pm r)$; the radial direction is the flat one (the scalon) and the gauge orbit has rank three. Every term is homogeneous of degree four, so each mass squared is proportional to $N$. For $k_{4} = k_{5} = 0$: $k_{8}\operatorname{Re}Q^{2} + k_{9}\operatorname{Im}Q^{2} \ge -r|Q|^{2}$, so $V \ge V_{4}(N, m) + (k_{10} - r)|Q|^{2}$ with $V_{4} = aN^{2} + bm^{2} + eNm$; on all of $\mathcal{R}_{\C}$, $N^{2} \ge |Q|^{2} + m^{2} \ge m^{2}$ (\claimref{synthesis.rung-field-on-the-lift-locus}(5)), and $V_{4} \ge 0$ for $|m| \le N$ under the admissibility conditions (\claimref{synthesis.scale-free-rung-potentials-need-the-flat-physical-branch}(3)); so $V \ge 0$ when $k_{10} \ge r$, the ray is a minimum of $V$ on the unit sphere as the imported Fact requires, and $m_{H,A}^{2} \ge m_{\pm}^{2}$. If $k_{10} + r < -(2a - e)/2$ one normal mass squared is negative (the check exhibits $V = -\tfrac{1}{200}$ next to the ray for $k_{3} = \tfrac12$, $k_{10} = -1$).

\emph{(4).} Fact \claimref{synthesis.imported-coleman-weinberg-gildener-weinberg-scalon-mass} applies to real scalars with a scale-free quartic polynomial potential and a flat direction at a minimum on the unit sphere; the linear field of (3) is such a theory (eight real scalars), and the cone $L_{0}$ is not a linear space, so the constrained field is outside the Fact's hypotheses except as a limit. Counting real scalars with masses at $\varphi = v$: the three gauge directions are massless in Landau gauge and the scalon is the flat direction, leaving the charged pair (two) and the normal pair (two); vectors and fermions are as in \claimref{synthesis.transmutation-of-the-rung-field-needs-a-charged-pair-near-450-gev}(i), under the same hypothesis (D). This gives the displayed relation, the standard sum rule of the scale-invariant two-doublet model: J. S. Lee and A. Pilaftsis, Phys. Rev. D 86 (2012) 035004 (arXiv:1201.4891), and K. Lane and W. Shepherd, arXiv:1808.07927, Eq. (10), $(\sum M^{4})^{1/4} = 540$ GeV over the heavy states (Fact \claimref{synthesis.imported-two-doublet-scalon-sum-rule-and-lep-neutrino-count}(i); their fermion couplings differ from the series' readings, which the sum rule does not use). Solving for $X = 2m_{\pm}^{4} + m_{H}^{4} + m_{A}^{4}$, the right side is $8\pi^{2}v^{2}m_{h}^{2} - 6M_{W}^{4} - 3M_{Z}^{4} + 12m_{t}^{4} + \ell$, the expression the stored result calls $2m_{\pm}^{4}$; it increases in $\pi$, $v$, $m_{h}$, $m_{t}$, and the stored check's corner evaluation, repeated here with $\pi \in [\tfrac{333}{106}, \tfrac{355}{113}]$, gives $2 \cdot 450^{4} < X < 2 \cdot 456^{4}$. Then $2m_{\pm}^{4} \le X$ gives $m_{\pm} < 456$; $m_{H}^{4} \le X$ and $X < 542^{4}$ give $m_{H}, m_{A} < 542$ (computed: $X^{1/4}$ lies between $535.8$ and $541.6$ GeV over the box); and $m_{H}, m_{A} \ge m_{\pm}$ give $4m_{\pm}^{4} \le X < 4 \cdot 383^{4}$. Equality $X = 2m_{\pm}^{4}$, the stored bracket, holds only for $m_{H} = m_{A} = 0$. The normal pair is the neutral component of the doublet $\alpha$, so the $Z$ couples to it with doublet strength and $Z \to HA$ is open when $m_{H} + m_{A} < M_{Z}$; a massless pair would add about half a neutrino species to the $Z$'s invisible width, which the measured $N_{\nu} = 2.9840 \pm 0.0082$ excludes (argued; Fact \claimref{synthesis.imported-two-doublet-scalon-sum-rule-and-lep-neutrino-count}(ii)); this does not change the bounds above appreciably, since $M_{Z}^{4}/8 \ll X$.

\emph{Reading.} In the scale-free case the bound $m_{H}^{2} < 542^{2}$ GeV$^{2}$ gives $4N(k_{10} + r) + m_{\pm}^{2} < 542^{2}$ GeV$^{2}$, so $k_{10} + r < 542^{2}/(2v^{2}) \approx 2.4$ (computed): the null-lift constraint is not imposed by a large coupling, and the elementary rung field under transmutation is two complete doublets with a neutral pair below $542$ GeV. With a level ($\xi \ne 0$, not scale free) $k_{10}$ is free and the constrained field is a limit of the linear one, but $k_{10} \gg 4\pi$ would put the normal pair's quartic beyond perturbation theory (argued). Transport-convenient throughout: the field, the locus, the unit and the boost are the series'; the spectrum and the scalon sum are standard two-doublet and Gildener-Weinberg computations.

\emph{(5).} With $f(\tau) = \arcsin^{2}\sqrt{\tau} = \tfrac12\sum_{n \ge 1}(4\tau)^{n}/(n^{2}\binom{2n}{n})$, whose $n = 1$ term is $\tau$, $A_{0} = (f - \tau)/\tau^{2} = \tfrac12\sum_{n \ge 2}4^{n}\tau^{n - 2}/(n^{2}\binom{2n}{n})$, a series with positive coefficients and leading term $\tfrac13$, increasing on $0 < \tau < 1$. On the flat ray $m_{\pm}^{2} \propto \varphi^{2}$, so the charged pair's term has unit weight, as in the stored result; the normal pair carries no electric charge. The Standard Model amplitude $A_{\mathrm{SM}} = A_{1}(\tau_{W}) + \tfrac43 A_{1/2}(\tau_{t})$ is negative with $|A_{\mathrm{SM}}| > 1 > A_{0}$ (the check brackets it below $-6.4$), so $\mu_{\gamma\gamma} = ((A_{\mathrm{SM}} + A_{0})/A_{\mathrm{SM}})^{2}$ decreases as $A_{0}$ grows, that is as $m_{\pm}$ decreases. At $m_{\pm} = 456$ GeV, with the stored $m_{h}$, $M_{W}$ and pole $m_{t}$, exact interval arithmetic gives $\mu_{\gamma\gamma} < 0.900$; so the bound holds for every $m_{\pm} \in (m_{h}/2, 456]$ GeV. The check also brackets the rate below $0.900$ and decreasing at $400$, $300$, $200$, $100$ and $70$ GeV (computed: about $0.899$, $0.898$, $0.895$, $0.873$, $0.806$; the heavy limit $A_{0} = \tfrac13$ gives $0.900$). Against the Particle Data Group average $1.10 \pm 0.06$ quoted in the stored result, the deficit is about three standard deviations for every charged mass, so freeing the charged mass through the normal pair does not relieve it.
